\begin{rubric}{学术概况}
\entry*[]{教育背景：上海交通大学生物医学工程博士（2024），现为波士顿学院博士后研究员。}
\entry*[]{研究方向：聚焦复杂三维生物结构的医学与电镜（EM）图像分析，重点研究弱/少监督条件下的三维实例分割方法，以及全细胞器分割基准与评测体系的构建。}
\entry*[]{方法特色：聚焦少量标注数据学习、一致性正则化、对比/对抗学习，覆盖 EM/CT 等多模态数据，强调“方法-数据集-评测”完整链条。}
\entry*[]{代表性成果：SCI期刊论文5篇（其中JCR Q1：2篇，第一/共同一作4篇），MICCAI/MLMI 等国际会议论文4篇。}
\entry*[]{基准与挑战：牵头 MitoEM 2.0 3D 线粒体实例分割基准（Scientific Data 在审）；CellMap Segmentation Challenge（289个eFIB-SEM数据、40+ 细胞器类别）第一名；MICCAI FLARE 2023（4000+ 多中心腹部CT，13器官+泛癌病灶分割）第二名，突出精度与效率。}
\entry*[]{在研工作：CR-GLoCo（CMIG 小修）及全细胞器分割、mitoVerse/SkelAffinity-Seg 等在研研究，计划投稿 IEEE TMI。}
\end{rubric}
